\chapter{Sea Ice Thermodynamics}\label{chap:ch11_seaice}

% --- local symbols (minimal, chapter-scoped) -----------------------
\providecommand{\hice}{h}              % ice thickness
\providecommand{\Aice}{A}              % ice concentration
\providecommand{\Tsi}{T_{s}}          % ice surface temperature
\providecommand{\Tfr}{T_{f}}           % freezing point of seawater
\providecommand{\Lf}{L_{f}}            % latent heat of fusion
\providecommand{\kice}{k_{i}}          % thermal conductivity of ice
\providecommand{\rhoice}{\rho_{i}}     % ice density
\providecommand{\Fnet}{F_{\mathrm{net}}}

Sea ice is the climate system's most leveraged component per unit mass.
A thin film of frozen seawater---typically $1$--$3\,\mathrm{m}$ thick---sits
between a cold polar atmosphere and an ocean held at its freezing point, and it
controls the heat, freshwater, and momentum exchanged across that interface. Two
feedbacks make it dynamically central. First, the \emph{ice--albedo feedback}:
open water absorbs roughly $90\%$ of incident shortwave radiation, bare ice
roughly $30\%$, so a small retreat of the ice edge warms the surface and
promotes further retreat. Second, the \emph{insulation feedback}: ice is a poor
conductor, so a growing ice cover throttles the ocean-to-atmosphere heat loss
that drives convection and, ultimately, the deep-water formation discussed in
Chapter~\ref{chap:ch08_ocean_overturning}. Sea ice also rejects brine as it
freezes and releases freshwater as it melts, coupling directly to the salinity
budget and hence to the overturning circulation studied in
Chapter~\ref{chap:ch17_amoc_tipping}.

\model{} represents these processes with a \emph{Semtner zero-layer}
thermodynamic model \citep{semtner1976}: the simplest scheme that captures the
surface energy balance, conductive growth, and the albedo/insulation feedbacks
without resolving the internal temperature profile of the ice or its dynamics
(drift, ridging, rheology). This chapter derives the governing equations, then
shows exactly which discrete equations the code integrates, where each
approximation, clamp, and limiter lives, and where the scheme departs from
observations. The two source modules are
\codref{chronos\_esm/ice/thermodynamics.py}{surface energy balance \& albedo}
and \codref{chronos\_esm/ice/driver.py}{growth/melt, frazil, fluxes, state}.

\begin{remark}
``Zero-layer'' means the ice slab has \emph{no} prognostic heat capacity: heat
conducts through it instantaneously and linearly. The full Semtner 1976 model
also offers three-layer and zero-layer variants with a snow layer and brine heat
storage; \model{} implements the bare zero-layer core (no snow, no internal heat
reservoir). This is honest to keep in mind: the model cannot represent the
thermal-inertia lag of thick multiyear ice.
\end{remark}

%====================================================================
\section{State variables and the freezing point}
%====================================================================

The prognostic state per ocean grid cell is the triple
\codref{chronos\_esm/ice/driver.py}{\code{IceState}}
\begin{equation}
  \bigl(\hice,\ \Aice,\ \Tsi\bigr),
\end{equation}
the grid-cell mean ice \emph{thickness} $\hice\,[\mathrm{m}]$, the fractional
ice \emph{concentration} $\Aice\in[0,1]$, and the ice \emph{surface
temperature} $\Tsi\,[^\circ\mathrm{C}]$. The model is initialised ice-free
($\hice=\Aice=0$, $\Tsi=\Tfr$) and forms ice dynamically.

Seawater freezes below $0^\circ\mathrm{C}$ because dissolved salt depresses the
freezing point. \model{} uses a constant
\begin{equation}
  \Tfr = -1.8\,^\circ\mathrm{C},
  \label{eq:tfreeze}
\end{equation}
the standard value for ocean salinity near $34$--$35\,\mathrm{psu}$. This is an
approximation: the true freezing point is salinity-dependent,
$\Tfr \approx -0.054\,\Salt$, so Eq.~\eqref{eq:tfreeze} is exact only near
$33\,\mathrm{psu}$ and biased by up to $\pm 0.1\,\mathrm{K}$ across the realistic
salinity range---negligible for the energy budget but worth noting for the
brackish Baltic or the salty subtropics. The constant is \code{T\_FREEZE} in
\code{thermodynamics.py}.

%====================================================================
\section{Surface energy balance}
%====================================================================

Over an ice surface the skin temperature $\Tsi$ adjusts so that the net heat
flux into the surface vanishes (an infinitesimally thin skin cannot store
energy). The continuous balance, all fluxes positive \emph{into} the surface,
is
\begin{equation}
  \Fnet(\Tsi) =
    \underbrace{(1-\alpha)\,S^{\downarrow}}_{\text{absorbed SW}}
    + \underbrace{L^{\downarrow}}_{\text{down LW}}
    - \underbrace{\varepsilon\sigma (\Tsi+T_0)^4}_{\text{up LW}}
    + \underbrace{\frac{\kice}{\hice}\,(\Tfr - \Tsi)}_{\text{conduction}}
    + \underbrace{C_b\,(T_a - \Tsi)}_{\text{turbulent}}
  = 0,
  \label{eq:ebalance}
\end{equation}
where $\alpha$ is the surface albedo, $S^{\downarrow}$ and $L^{\downarrow}$ are
the downward shortwave and longwave fluxes $[\mathrm{W\,m^{-2}}]$,
$\varepsilon=0.97$ is the emissivity, $\sigma=5.67\times10^{-8}\,
\mathrm{W\,m^{-2}\,K^{-4}}$ the Stefan--Boltzmann constant,
$T_0=273.15\,\mathrm{K}$, $T_a$ the air temperature, and $\kice=2.03\,
\mathrm{W\,m^{-1}\,K^{-1}}$ the conductivity of ice. The conduction term is the
zero-layer heart of the model: it linearises the temperature profile across the
slab between the freezing ocean at the base ($\Tfr$) and the radiating skin
($\Tsi$), so the conductive flux is simply $\kice\,(\Tfr-\Tsi)/\hice$.

The turbulent (sensible $+$ latent) fluxes are bulk-linearised to
$C_b(T_a-\Tsi)$ with a lumped exchange coefficient $C_b = 10\,
\mathrm{W\,m^{-2}\,K^{-1}}$. This is a deliberate simplification: it absorbs the
sensible and latent heat into one temperature-difference term and omits the
wind-speed and humidity dependence a full bulk formula carries (contrast the
coupler bulk fluxes in Chapter~\ref{chap:ch13_coupler_fluxes}).

\paragraph{Albedo.} The albedo is not a step function but a doubly-smooth
sigmoid blend, chosen for differentiability (see
Chapter~\ref{chap:ch03_differentiable}). Writing
$s(x)=(1+e^{-x})^{-1}$ for the logistic function,
\codref{chronos\_esm/ice/thermodynamics.py}{\code{compute\_albedo}}
\begin{align}
  m  &= s\!\left(\frac{\Tsi+0.5}{0.5}\right),
      &\alpha_{\mathrm{ice}} &= (1-m)\,\alpha_{i} + m\,\alpha_{m}, \\
  f_{\mathrm{thin}} &= s\!\left(\frac{0.5-\hice}{0.1}\right),
      &\alpha &= (1-f_{\mathrm{thin}})\,\alpha_{\mathrm{ice}}
                 + f_{\mathrm{thin}}\,\alpha_{\mathrm{ocean}},
  \label{eq:albedo}
\end{align}
with $\alpha_i=0.7$ (cold bare ice), $\alpha_m=0.5$ (melting ice/melt ponds),
and $\alpha_{\mathrm{ocean}}=0.1$ (open water). The first sigmoid melts the
albedo down as $\Tsi \to 0^\circ\mathrm{C}$ over a $0.5\,\mathrm{K}$ width; the
second blends toward the ocean value for thin ice ($\hice \lesssim 0.5\,
\mathrm{m}$), so a thin, broken cover is not given the full reflectivity of a
thick floe. Both transitions are smooth and therefore differentiable---a hard
threshold would zero the gradient $\partial\alpha/\partial\hice$ used by the
adjoint applications of Chapter~\ref{chap:ch15_differentiable_applications}.

\paragraph{Solving for $\Tsi$.} Equation~\eqref{eq:ebalance} is nonlinear in
$\Tsi$ through the quartic longwave term, so the code takes a fixed number of
Newton steps---three---rather than iterating to a tolerance. A fixed step count
keeps the solver a static, differentiable computational graph (\jax{} can trace
and back-propagate through it). With
$\Fnet'(\Tsi) = -4\varepsilon\sigma(\Tsi+T_0)^3 - \kice/\hice - C_b$,
\codref{chronos\_esm/ice/thermodynamics.py}{\code{solve\_surface\_temp}}
\begin{equation}
  \Tsi^{(n+1)} = \Tsi^{(n)}
      - \frac{\Fnet(\Tsi^{(n)})}{\Fnet'(\Tsi^{(n)}) - 10^{-5}},
  \qquad n=0,1,2,
  \label{eq:newton}
\end{equation}
seeded with $\Tsi^{(0)}=T_a$. The $-10^{-5}$ guard keeps the denominator away
from zero. The thickness entering $\kice/\hice$ is floored at
$\hice_{\mathrm{safe}}=\max(\hice,\,0.1\,\mathrm{m})$ to avoid the singular
conduction of vanishing ice. Finally the skin temperature is capped at the
melting point,
\begin{equation}
  \Tsi \leftarrow \min(\Tsi,\ 0\,^\circ\mathrm{C}),
  \label{eq:tcap}
\end{equation}
and $\Fnet$ is re-evaluated at the capped $\Tsi$. This cap is the mechanism
that distinguishes growth from melt: if the radiative balance would drive
$\Tsi>0$, the surface cannot warm past melting, and the leftover positive
$\Fnet$ becomes melt energy (Section~\ref{sec:growth}).

\begin{remark}
Three Newton steps from $\Tsi^{(0)}=T_a$ is usually enough because the air
temperature is a good first guess and the quartic is gently curved over the
relevant range. It is \emph{not} guaranteed to converge under extreme forcing;
the melt cap~\eqref{eq:tcap} and the downstream clamps
(Section~\ref{sec:growth}) provide the safety net that an exact solve would
otherwise give. This trades a little physical fidelity for a fixed, fast,
differentiable graph.
\end{remark}

%====================================================================
\section{Ice growth and melt}\label{sec:growth}
%====================================================================

The slab thickness changes through three pathways, all phrased as a Stefan
condition---heat exchanged at a phase boundary converts to a thickness rate via
the latent heat of fusion $\Lf = 3.34\times10^5\,\mathrm{J\,kg^{-1}}$ and ice
density $\rhoice = 917\,\mathrm{kg\,m^{-3}}$:
\begin{equation}
  \dt{\hice} = \underbrace{\frac{\kice(\Tfr-\Tsi)/\hice_{\mathrm{safe}}}
                                {\rhoice \Lf}}_{\text{bottom growth}}
             - \underbrace{\frac{Q_{o}}{\rhoice \Lf}}_{\text{bottom melt}}
             - \underbrace{\frac{\max(\Fnet,0)}{\rhoice \Lf}}_{\text{top melt}}.
  \label{eq:dhdt}
\end{equation}
Each term has a clear physical reading.

\paragraph{Bottom (congelation) growth.} When the surface is cold ($\Tsi<0$)
the energy balance is satisfied and $\Fnet\approx 0$; the ice grows from below
because the conductive flux $\kice(\Tfr-\Tsi)/\hice$ carries the ocean's latent
heat away through the slab to the cold sky. Thicker ice conducts less, so growth
self-limits as $\hice^{-1}$---the classic $\sqrt{t}$ thickening law of Stefan.

\paragraph{Top melt.} When the radiative balance would push $\Tsi$ above
$0^\circ\mathrm{C}$, the cap~\eqref{eq:tcap} holds it at melting and the residual
$\max(\Fnet,0)>0$ is spent melting ice from the top.

\paragraph{Bottom melt by ocean heat.} If the ocean mixed layer is above
freezing, it delivers heat to the ice base. \model{} parameterises this ocean
heat flux as a relaxation of the mixed layer's excess temperature toward
$\Tfr$,
\codref{chronos\_esm/ice/driver.py}{ocean heat flux, \code{step\_ice}}
\begin{equation}
  Q_{o} = \frac{\rho_w c_w\,\max(T_{\mathrm{sst}}-\Tfr,0)\,H_{\mathrm{ml}}}
               {\tau_{\mathrm{mix}}},
  \label{eq:ohf}
\end{equation}
with $\rho_w=1025\,\mathrm{kg\,m^{-3}}$, $c_w=3985\,\mathrm{J\,kg^{-1}\,K^{-1}}$,
mixed-layer depth $H_{\mathrm{ml}}=50\,\mathrm{m}$, and exchange timescale
$\tau_{\mathrm{mix}}=10\,\mathrm{days}$. This yields $O(10\,
\mathrm{W\,m^{-2}})$ for a few tenths of a degree of excess heat, in the observed
polar range, but the fixed $H_{\mathrm{ml}}$ and $\tau_{\mathrm{mix}}$ are
crude---there is no prognostic mixed layer here.

\paragraph{Discrete update.} The driver integrates Eq.~\eqref{eq:dhdt} with one
explicit Euler step at the ocean timestep $\Delta t = 900\,\mathrm{s}$
(\code{DT\_OCEAN}):
\begin{equation}
  \hice^{n+1} = \hice^{n} + \Delta t\,\dt{\hice}.
\end{equation}

%====================================================================
\section{Frazil: new-ice formation}
%====================================================================

Where there is no ice yet but the ocean is supercooled
($T_{\mathrm{sst}}<\Tfr$), \model{} forms new ice (\emph{frazil}) by converting
the heat deficit of the mixed layer directly into thickness. The supercooling is
clamped to $2\,\mathrm{K}$ before conversion, then the new ice is rate-limited to
$0.1\,\mathrm{m}$ per step,
\codref{chronos\_esm/ice/driver.py}{frazil block, \code{step\_ice}}
\begin{align}
  \Delta T_{sc} &= \min\!\bigl(\max(\Tfr - T_{\mathrm{sst}},0),\ 2\,\mathrm{K}\bigr),\\
  \hice_{\mathrm{new}} &= \min\!\left(
        \frac{\rho_w c_w\,\Delta T_{sc}\,H_{\mathrm{ml}}}{\rhoice \Lf},\
        0.1\,\mathrm{m}\right),
  \label{eq:frazil}
\end{align}
which is added to $\hice^{n+1}$. The two limiters are numerical hygiene, not
physics: without them a single cold step over a $50\,\mathrm{m}$ mixed layer
could dump several metres of ice instantaneously (a $5\,\mathrm{K}$ deficit
$\times\,50\,\mathrm{m}$ is $\sim 6\,\mathrm{m}$ of ice), destabilising the
coupled run. The comments in the source record that the supercooling clamp was
tightened from $5\,\mathrm{K}$ to $2\,\mathrm{K}$ for exactly this reason. The
trade-off is honest: real frazil in a strongly supercooled lead would form
faster than $0.1\,\mathrm{m}/900\,\mathrm{s}$, so the model lags rapid opening
events.

After growth, melt, and frazil, the thickness is clamped to a physically
realistic ceiling and floor,
\begin{equation}
  \hice^{n+1} \leftarrow \max\!\bigl(\min(\hice^{n+1},\ 5\,\mathrm{m}),\ 0\bigr),
  \label{eq:hclamp}
\end{equation}
and masked to ocean cells ($\hice \leftarrow \hice\cdot\text{mask}$, land $=0$).
The $5\,\mathrm{m}$ cap is a guard against an earlier failure mode in which
unbounded conductive growth produced $\sim 39\,\mathrm{m}$ ice; $5\,\mathrm{m}$
is at the upper end of multiyear Arctic ice. Capping thickness this way
suppresses pressure ridges and the thickest deformed ice, a known low bias of
the scheme.

%====================================================================
\section{Concentration}
%====================================================================

The zero-layer model carries no ice dynamics, so concentration is not advected
or ridged; it is diagnosed from thickness by a smooth saturating map,
\codref{chronos\_esm/ice/driver.py}{concentration, \code{step\_ice}}
\begin{equation}
  \Aice = \tanh\!\left(\frac{\hice}{0.5\,\mathrm{m}}\right).
  \label{eq:conc}
\end{equation}
A cell with $\hice\gtrsim 1\,\mathrm{m}$ is essentially fully covered
($\Aice\to1$), while thin ice gives partial cover. This is a Hibler-style
parameterisation in spirit but far simpler: there is no separate open-water
(lead) fraction equation and no thermodynamic redistribution. Because $\Aice$ is
a monotone function of $\hice$, the model cannot represent thick, low-concentration
fields (broken pack) or thin, fully consolidated new ice. Treat $\Aice$ here as a
diagnostic of how ``complete'' the cover is, not an independent dynamical
variable.

%====================================================================
\section{How ice modifies the ocean--atmosphere fluxes}
%====================================================================

The reason a thermodynamic ice model matters to the rest of \model{} is the heat
and freshwater it hands to the ocean. \code{step\_ice} returns the pair
$(Q_{\mathrm{ocn}},\,F_{\mathrm{fw}})$.

\paragraph{Heat flux to the ocean.} Two ice processes exchange heat with the
ocean below:
\codref{chronos\_esm/ice/driver.py}{flux block, \code{step\_ice}}
\begin{equation}
  Q_{\mathrm{ocn}} = \underbrace{-\,\frac{\kice(\Tfr-\Tsi)}{\hice_{\mathrm{safe}}}}
                        _{\text{conductive extraction}}
                   + \underbrace{\frac{\hice_{\mathrm{new}}\,\rhoice \Lf}{\Delta t}}
                        _{\text{frazil latent release}}.
  \label{eq:qocn}
\end{equation}
Congelation growth \emph{removes} heat from the ocean--ice interface (the first,
negative term: the ocean must give up latent heat to freeze), whereas frazil
formation \emph{releases} latent heat back to the mixed layer as the supercooled
water freezes (the second, positive term). The sign convention is that
$Q_{\mathrm{ocn}}$ is energy delivered to the ocean. Crucially, while ice is
present the ocean is insulated from the atmospheric radiative and turbulent
fluxes---that masking is applied in the coupler
(Chapter~\ref{chap:ch13_coupler_fluxes}), and Eq.~\eqref{eq:qocn} is the
\emph{only} thermal channel left open between ocean and the surface there. This
is the insulation feedback in code form.

\paragraph{Freshwater flux.} Ice growth concentrates the seawater it leaves
behind (brine rejection raises ocean salinity); melt does the reverse. The model
expresses this as a freshwater flux proportional to the thickness tendency,
\begin{equation}
  F_{\mathrm{fw}} = -\,\rhoice\,\frac{\hice^{n+1}-\hice^{n}}{\Delta t}
  \qquad [\mathrm{kg\,m^{-2}\,s^{-1}}].
  \label{eq:fw}
\end{equation}
Growth ($\hice$ increasing) gives $F_{\mathrm{fw}}<0$: freshwater is locked into
the ice, leaving the ocean saltier. Melt ($\hice$ decreasing) gives
$F_{\mathrm{fw}}>0$: meltwater freshens the surface. This coupling to the
salinity field is precisely the high-latitude freshwater forcing that the
overturning circulation is sensitive to
(Chapters~\ref{chap:ch08_ocean_overturning} and~\ref{chap:ch17_amoc_tipping}).

\begin{remark}[Honest accounting note]
Equation~\eqref{eq:fw} converts \emph{all} of the net thickness change---including
the frazil term~\eqref{eq:frazil}---into a freshwater flux, while the heat
flux~\eqref{eq:qocn} accounts for frazil latent heat separately. Because frazil
freezing also rejects brine, treating its mass as freshwater here is consistent;
but note that the scheme uses pure-ice density $\rhoice$ and assumes the ice is
salt-free (zero ice salinity). Real sea ice retains $\sim 4$--$8\,\mathrm{psu}$
of brine, so the model slightly overstates the freshwater (and brine-rejection)
signal per metre of ice. The bias is second-order for the energy budget but
relevant for precise salinity tendencies.
\end{remark}

%====================================================================
\section{Summary of the time step}
%====================================================================

Per ocean step $\Delta t$, \code{step\_ice} executes, in order: (1) solve the
surface energy balance for $\Tsi$ by three Newton iterations and cap at melting
[Eqs.~\eqref{eq:ebalance}--\eqref{eq:tcap}]; (2) form the bottom-growth,
top-melt, and ocean-heat-melt terms and Euler-step the thickness
[Eq.~\eqref{eq:dhdt}]; (3) add rate-limited frazil where the ocean is supercooled
[Eq.~\eqref{eq:frazil}]; (4) clamp thickness to $[0,5]\,\mathrm{m}$ and mask land
[Eq.~\eqref{eq:hclamp}]; (5) diagnose concentration [Eq.~\eqref{eq:conc}]; and
(6) return the ocean heat and freshwater fluxes
[Eqs.~\eqref{eq:qocn}--\eqref{eq:fw}]. The scheme is deliberately minimal: it
omits snow, internal heat storage, ice dynamics/advection, ridging, a prognostic
lead fraction, and a salinity-dependent freezing point. Within those limits it
faithfully reproduces the two feedbacks that make sea ice climatically
important---albedo and insulation---and remains end-to-end differentiable.

\begin{table}[t]
\centering
\caption{Physical constants in the \model{} sea-ice thermodynamics
(\code{thermodynamics.py}, \code{driver.py}).}
\begin{tabular}{lll}
\toprule
Symbol & Value & Meaning \\
\midrule
$\rhoice$            & $917\,\mathrm{kg\,m^{-3}}$              & ice density \\
$\Lf$                & $3.34\times10^{5}\,\mathrm{J\,kg^{-1}}$ & latent heat of fusion \\
$\kice$              & $2.03\,\mathrm{W\,m^{-1}\,K^{-1}}$      & ice thermal conductivity \\
$\Tfr$               & $-1.8\,^\circ\mathrm{C}$                & freezing point of seawater \\
$\varepsilon$        & $0.97$                                  & ice emissivity \\
$C_b$                & $10\,\mathrm{W\,m^{-2}\,K^{-1}}$        & bulk turbulent coefficient \\
$\alpha_i,\alpha_m,\alpha_{\mathrm{ocean}}$ & $0.7,\,0.5,\,0.1$ & cold/melt/ocean albedo \\
$\tau_{\mathrm{mix}}$ & $10\,\mathrm{days}$                    & ice--ocean exchange timescale \\
$H_{\mathrm{ml}}$    & $50\,\mathrm{m}$                        & assumed mixed-layer depth \\
\bottomrule
\end{tabular}
\end{table}

%====================================================================
\section*{Exercises}
\addcontentsline{toc}{section}{Exercises}
%====================================================================

\begin{exercise}[The Stefan thickening law]
Set the ocean heat flux $Q_o=0$ and the surface fluxes so that $\Fnet=0$
(perpetual polar night, $\Tsi$ fixed at some $T_s<\Tfr$ by strong longwave
cooling). Equation~\eqref{eq:dhdt} then reduces to
$\dot{\hice}=\kice(\Tfr-T_s)/(\rhoice\Lf\,\hice)$. Integrate analytically to show
$\hice(t)=\sqrt{2\kice(\Tfr-T_s)t/(\rhoice\Lf)}$. Now run \code{step\_ice} in a
single-column driver with constant forcing and $\hice^0=0.1\,\mathrm{m}$,
$T_s=-20^\circ\mathrm{C}$, and compare the simulated $\hice(t)$ to the
$\sqrt{t}$ law. At what thickness does the $5\,\mathrm{m}$ clamp
[Eq.~\eqref{eq:hclamp}] start to distort the curve, and how many days does that
take?
\end{exercise}

\begin{exercise}[Albedo feedback and differentiability]
Using Eq.~\eqref{eq:albedo}, plot $\alpha$ as a function of $\hice$ at fixed
$\Tsi=-5^\circ\mathrm{C}$ and as a function of $\Tsi$ at fixed
$\hice=2\,\mathrm{m}$. Then use \jax{}'s \code{jax.grad} on
\code{compute\_albedo} to obtain $\partial\alpha/\partial\hice$ and
$\partial\alpha/\partial\Tsi$. Where are the gradients largest, and why would a
hard ($\tanh\!\to\!$ step) threshold break the adjoint applications of
Chapter~\ref{chap:ch15_differentiable_applications}?
\end{exercise}

\begin{exercise}[The insulation feedback in the coupled model]
Insulation enters through the conduction term $\kice/\hice$. Run a short coupled
simulation, then artificially \emph{double} $\kice$ in \code{thermodynamics.py}
and rerun from the same initial state. Predict, then verify, the sign of the
change in (i) equilibrium ice thickness, (ii) the conductive heat flux to the
ocean $Q_{\mathrm{ocn}}$ [Eq.~\eqref{eq:qocn}], and (iii) high-latitude ocean
convection. Explain why more conductive ice can paradoxically be \emph{thinner}
ice.
\end{exercise}

\begin{exercise}[Brine rejection and the overturning]
The freshwater flux~\eqref{eq:fw} ties sea-ice growth to ocean salinity. Estimate
the brine flux from freezing $1\,\mathrm{m}$ of ice over a winter in a
$3.75^\circ$ T31 polar cell, expressing it as an equivalent salinity tendency for
a $50\,\mathrm{m}$ mixed layer. Then discuss qualitatively how the zero-ice-salinity
assumption (Remark, Section~6) biases this estimate, and connect it to the
high-latitude freshwater sensitivity of the AMOC in
Chapter~\ref{chap:ch17_amoc_tipping}.
\end{exercise}
